\newcommand{\NearTau}{\mathord{\text{\fontspec{Times New Roman}τ}}}%
\expandafter\def\csname HistoricalNoteText6\endcsname{Concessa explicatione phaenomenorum lucis in systemate emanationis, refractio pendet ab attractione particularum corporis pellucidi in particulas lucis moleculari, ratioque refractionis a valore ipsius $\psi0$, ita quidem ut habeatur
\[
\frac{c\psi0}{g}=\frac{(nn-1)kk}{8\pi^3l}
\]
denotante $l$ longitudinem penduli per minutum secundum vibrantis, $k$ motum luminis in vacuo intra minutum secundum, $n$ rationem sinus anguli incidentiae ad sinum anguli refracti: hoc pacto pro aqua fit $\frac{c\psi0}{g}$ bis millies maior quam distantia media solis a terra.}%
\SourceLeaf{52}terum manifesto $\frac{c\psi r}{r}$ est quantitas cum gravitate $g$ homogenea, sive $\frac{c\psi r}{g}$ linea, adeoque $\frac{c\psi0}{g}$ linea determinata (pro natura corporum, ad quorum vires attractivas functio $fr$ refertur), cuius magnitudinem ingentem suspicari quidem licet, sed quam in casibus determinatis vix approximative assignare valemus, saltem non absque suppositionibus hypotheticis\HistoricalNote{6}.

\Article{13}
Prorsus simili modo in integratione $\int\psi r.\mathrm{d}r=-\theta r$ constantem ita electam supponimus, ut fiat $\theta r=0$ pro valore arbitrario ipsius $r$ intra ambitum eorum, pro quibus experimenta instituere licet, quo pacto $\theta r$ insensibilis erit pro quovis eiusmodi valore sensibili ipsius $r$, etiamsi sensibilis evadere possit pro valore insensibili. Manifesto $\frac{c\theta r}{g}$ exprimit aream figurae duarum dimensionum, adeoque $\frac{\theta r}{\psi r}$ lineam. Necessario autem $\frac{\theta0}{\psi0}$ est linea magnitudinis insensibilis, quod ita demonstramus. Quum $\psi r$ inde a $r=0$ continuo decrescat, et quidem tam cito, ut iam insensibilis evaserit, quamprimum $r$ valorem sensibilem acquisivit, valor ipsius $r$, pro quo fit $\psi r=\frac12\psi0$, insensibilis esse debet: denotetur ille per $\rho$. Consideremus integrale $\int(\psi0-\psi r)\mathrm{d}r$, quod ab $r=0$ usque ad $r=R$ extensum fit $=R\psi0-\theta0+\theta R$. Manifesto hoc integrale maius erit, quam idem integrale ab $r=\rho$ usque ad $r=R$ extensum, atque hoc iterum maius, quam integrale $\int(\psi0-\psi\rho)\mathrm{d}r$ inter eosdem limites. Quare quum integrale postremum fiat $=(\psi0-\psi\rho)(R-\rho)=\frac12\psi0.(R-\rho)$, erit generaliter pro quovis valore ipsius $R$ (maiori quam $\rho$)
\[
R\psi0-\theta0+\theta R>\frac12\psi0.(R-\rho)
\]
Iam si $R$ denotare supponitur valorem fractionis $\frac{\theta0}{\psi0}$, haec relatio suppeditat
\[
\theta R>\frac12\psi0.(R-\rho)
\]
quod foret absurdum, si $R$ esset quantitas sensibilis.

\SourceLeaf{53}Non obstante itaque ingente magnitudine ipsius $\psi0$, nihil impedit, quominus $\theta0$ esse possit quantitas satis modica et cum dimensionibus corporum experimentis subiectorum comparabilis.

\Article{14}
Superest, ut quae ex hac indole functionis $\theta$ respectu integralis (I)
\[
\int\int\frac{\mathrm{d}t.\mathrm{d}T.\cos q.\cos Q.\theta(\mathrm{d}t,\mathrm{d}T)}{(\mathrm{d}t,\mathrm{d}T)^2}
\]
sequuntur, perscrutemur. Haec investigatio inchoare debet a simpliciori, dum in alterutra superficie punctum determinatum $\mu$ consideramus atque integrale (II)
\[
\int\frac{\mathrm{d}t.\cos q.\cos Q.\theta(\mu,\mathrm{d}t)}{(\mu,\mathrm{d}t)^2}
\]
per totam superficiem $t$ extendendum evolvimus. Denotant hic $Q$ angulum inter duas rectas a puncto $\mu$ proficiscentes, alteram versus elementum $\mathrm{d}t$, alteram fixam; $q$ vero angulum inter duas rectas a puncto elementi $\mathrm{d}t$ proficiscentes, alteram versus punctum $\mu$, alteram elemento normalem extrorsumque directam.

Primo loco observamus, si punctum $\mu$ sit in distantia sensibili a superficie $t$, valores omnium $\theta(\mu,\mathrm{d}t)$ insensibiles fore: in hoc itaque casu totum integrale (II) insensibile erit. Hoc itaque integrale eatenus tantum valorem sensibilem acquirere potest, quatenus superficies $t$ offert partes in distantia insensibili a puncto $\mu$ positas, manifestoque sufficit, integrale (II) per tales partes extendere, neglectis omnibus, quae sunt in distantiis sensibilibus.

Porro pro $\frac{\mathrm{d}t.\cos q}{(\mu,\mathrm{d}t)^2}$ restituemus $\pm\mathrm{d}\Pi$, denotante $\mathrm{d}\Pi$ in superficie sphaerica radio $=1$ circa centrum $\mu$ descripta elementum id, in quod elementum $\mathrm{d}t$ inde a puncto $\mu$ visum proiicitur, et valente signo superiori vel inferiori, prout elementum $\mathrm{d}t$ plagam exteriorem vel interiorem puncto $\mu$ advertit. Hoc pacto integrale (II) ita exhibetur
\[
\int\pm\mathrm{d}\Pi.\cos Q.\theta(\mu,\mathrm{d}t)
\]
patetque, huius integralis valorem eatenus tantum sensibilem fieri posse, quatenus elementa $\mathrm{d}\Pi$ talia, quae ad distantias insensibiles $(\mu,\mathrm{d}t)$ referuntur, spatium magnitudinis sensibilis in superficie sphaerica explent.

Hinc facile colligitur, integrale nostrum, generaliter loquendo, etiam insensibile manere, quoties punctum $\mu$ iaceat in ipsa superficie $t$: patet enim, pro\SourceLeaf{54}iectiones omnium elementorum $\mathrm{d}t$ a puncto $\mu$ insensibiliter remotorum esse in distantia insensibili a circulo maximo, quem format in superficie sphaerica planum superficiem $t$ in puncto $\mu$ tangens. Excipere oportet tres casus, puta

1) eum, ubi radii curvaturae superficiei $t$ in puncto $\mu$ sunt magnitudinis insensibilis.

2) eum, ubi continuitas curvaturae in puncto $\mu$, vel intra distantiam insensibilem ab eo interrumpitur (Conf.\ \textit{Disquiss. gen. circa superficies curvas} art.\ 3).

3) eum, ubi superficies $t$ offert partem aliam a puncto $\mu$ insensibiliter distantem, puta si apud hoc punctum crassities spatii $s$ est insensibilis. Ceterum huncce casum ei, quem in art.\ seq.\ tractabimus, adnumerare licet.

\Article{15}
Superest scilicet casus, ubi punctum $\mu$ non est in superficie ipsa $t$, attamen in distantia insensibili ab ea: in hoc casu integrale nostrum utique valorem sensibilem habere potest, quem iam accuratius examinabimus.

Secent superficiem sphaericam recta a puncto $\mu$ normaliter in superficiem $t$ ducta, atque recta fixa ibinde proficiscens resp.\ in punctis $G$, $H$; statuatur arcus $GH=k$, arcus autem inter $G$ atque punctum indefinitum superficiei sphaericae $=v$; denique sit $w$ angulus sphaericus inter arcus $k$, $v$. Hoc pacto pro elemento $\mathrm{d}\Pi$ accipere licet productum $\sin v.\mathrm{d}v.\mathrm{d}w$, unde scribendo brevitatis caussa $r$ pro $(\mu,\mathrm{d}t)$, integrale (II) fit
\[
=\int\int\pm(\cos k.\cos v+\sin k.\sin v.\cos w)\theta r.\sin v.\mathrm{d}v.\mathrm{d}w
\]
quam integrationem extendere tantummodo oportet per eas partes superficiei sphaericae, in quas distantiae insensibiles $r$ proiiciuntur. Referuntur hae ad partem insensibilem superficiei $t$, quam si pro \textit{plana} habemus, distantiamque minimam (puncto $G$ seu valori $v=0$ respondentem) per $\rho$ denotamus, fit $r=\frac{\rho}{\cos v}$, sive a $w$ independens; perfecta itaque integratione respectu variabilis $w$, puta a $w=0$ usque ad $w=360^0$, integrale fit
\[
=\pm\int2\pi\theta r.\cos k.\cos v.\sin v.\mathrm{d}v=\pm\int\frac{2\pi\cos k.\rho\rho\theta r.\mathrm{d}r}{r^3}
\]
quae integratio extendenda est ab $r=\rho$ usque ad valorem sensibilem arbitrarium quantumvis parvum. Statuendo itaque generaliter
\[
2rr\int\frac{\theta r.\mathrm{d}r}{r^3}=-\theta'r
\]
\SourceLeaf{55}accepta constante integrationis, ita ut fiat $\int\frac{\theta r.\mathrm{d}r}{r^3}=0$, pro valore arbitrario sensibili intra ambitum eorum, circa quos experimenta instituere licet, erit integrale (II), neglectis insensibilibus,
\[
=\pm\pi\cos k.\theta'\rho
\]

Si dubium videretur, utrum fas sit, partem superficiei $t$ intra distantiam insensibilem a puncto $\mu$ positam pro plana habere, consideremus eius loco sphaericam, et quidem sit $R$ distantia centri sphaerae a puncto $\mu$ positive vel negative sumenda, prout centrum est in directione versus $G$ vel in opposita. Ita erit
\[
\begin{gathered}
\cos v=\frac{\rho}{r}(1-\frac{\rho}{2R})+\frac{r}{2R}\\
\sin v.\mathrm{d}v=[\frac{\rho}{rr}(1-\frac{\rho}{2R})-\frac{1}{2R}]\mathrm{d}r
\end{gathered}
\]
unde facile colligitur, integrale pro hoc casu non differre quantitate sensibili a valore prius invento $\pm\pi\cos k.\theta'\rho$, si modo $R$ sit quantitas sensibilis. Quaecunque autem sit curvatura superficiei $t$ in ea parte, de qua agitur, dummodo radii curvaturae non sint insensibiles. semper duae superficies sphaericae assignari poterunt, superficiem $t$ in puncto ipsi $\mu$ proximo tangentes, intra quas $t$ sita sit, et quarum radii sint magnitudinis sensibilis, manifestoque tunc integrale nostrum intra integralia ad illas superficies relata cadet, et proin absque errore sensibili per eandem formulam exprimetur, quae tunc tantummodo exceptionem patitur, ubi superficies $t$ in distantia insensibili a puncto $\mu$ vel curvaturam radii insensibilis, vel aciem vel cuspidem offert.

\Article{16}
Quodsi iam ab integratione (II) ad integrale (I) progredimur, manifestum est, hoc insensibile fieri, non solum in eo casu, ubi illa pro \textit{nullo} puncto superficiei $T$ valorem sensibilem produxit, sed in eo quoque, ubi complexus elementorum superficiei $T$, pro quorum punctis integrale (II) sensibile evaserat, aream tantummodo insensibilis magnitudinis sistit. Quae si rite perpenduntur, apparebit, integrale (I) eatenus tantum valorem sensibilem acquirere posse, quatenus superficies $T$ partem vel partes sensibilis magnitudinis contineat in distantia insensibili a superficie $t$ positas. Quales partes quum a parallelismo cum superficie $t$ sensibiliter deviare nequeant, pro quovis earum puncto $\cos k$ non sensibi\SourceLeaf{56}liter differet vel a $+1$ vel a $-1$, prout plaga superficiei $T$ exterior vel interior superficiei $t$ advertitur. Quodsi itaque per $\NearTau$, $\NearTau'$ eas partes superficiei $T$ denotamus, quae sunt in distantia insensibili a superficie $t$, et quidem per $\NearTau$ eas, ubi plaga exterior alterius superficiei plagae interiori alterius advertitur, per $\NearTau'$ autem eas, ubi plagae homonymae sibi mutuo obvertuntur, denique per $\rho$ distantiam minimam cuiusvis elementi $\mathrm{d}\NearTau$ vel $\mathrm{d}\NearTau'$ a superficie $t$, integrale nostrum (I) neglectis insensibilibus fit
\[
=-\int\pi\theta'\rho.\mathrm{d}\NearTau+\int\pi\theta'\rho.\mathrm{d}\NearTau'
\]
Manifesto hic nihil interest. utrum partes $\NearTau$, $\NearTau'$ ad superficiem $T$ an ad $t$ referantur.

Hoc itaque modo iam nacti sumus solutionem completam problematis, quod in art.\ 6 nobis proposueramus, pro ea functionis $\varphi$ indole, cui tamquam basi disquisitio principalis de figura aequilibrii fluidorum innititur, scilicet habemus
\[
\int\int\mathrm{d}s.\mathrm{d}S.\varphi(\mathrm{d}s,\mathrm{d}S)=4\pi\sigma\psi0-\pi\tau\theta0+\pi\tau'\theta0-\pi\int\mathrm{d}\NearTau.\theta'\rho+\pi\int\mathrm{d}\NearTau'.\theta'\rho
\]

\Article{17}
Origo functionis $\theta'$ ita enunciari potest, ut sit
\[
\frac{\theta'r}{rr}=\int\frac{2\theta x.\mathrm{d}x}{x^3}
\]
sumto integrali ab $x=r$ usque ad valorem constantem sensibilem arbitrarium, quem hic per $R$ denotamus. Manifesto hoc integrale minus erit quam hoc $\int\frac{2\theta r.\mathrm{d}x}{x^3}$ inter eosdem limites, quod est $=\frac{\theta r}{rr}-\frac{\theta r}{RR}$, adeoque a potiori minus quam $\frac{\theta r}{rr}$. Quum autem indefinite habeatur
\[
\int\frac{2\theta x.\mathrm{d}x}{x^3}=-\frac{\theta x}{xx}+\int\frac{\mathrm{d}\theta x}{xx}=-\frac{\theta x}{xx}-\int\frac{\psi x.\mathrm{d}x}{xx}
\]
erit
\[
\frac{\theta'r}{rr}=\frac{\theta r}{rr}-\frac{\theta R}{RR}-\int\frac{\psi x.\mathrm{d}x}{xx}
\]
sumto integrali inter eosdem limites, quod minus erit quam integrale $\int\frac{\psi r\mathrm{d}x}{xx}$, adeoque etiam minus quam $\frac{\psi r}{r}$; quocirca valor ipsius $\frac{\theta'r}{rr}$ maior erit quam
\[
\frac{\theta r}{rr}-\frac{\theta R}{RR}-\frac{\psi r}{r}
\]
\SourceLeaf{57}Cadit itaque $\theta'r$ inter limites
\[
\theta r\quad\text{atque}\quad\theta r-rr.\frac{\theta R}{RR}-r\psi r
\]
quorum differentia, decrescente $r$ in infinitum, manifesto quavis quantitate assignabili minor evadere potest, quum supponamus esse vel $\psi0$ quantitatem finitam. Colligimus hinc, statui debere $\theta'0=\theta0$. Patet itaque, in formula, ad quam in art.\ praec.\ pervenimus, terminum $-\pi\tau\theta0$ tamquam sub termino $-\pi\int\mathrm{d}\NearTau.\theta'\rho$, atque terminum $\pi\tau'\theta0$ tamquam sub termino $\pi\int\mathrm{d}\NearTau'.\theta'\rho$ comprehensum considerari posse, si distinctionem, quam inter distantiam insensibilem et distantiam nullam fecimus, tolleremus, atque partes $\tau$, $\tau'$ resp.\ partibus $\NearTau$, $\NearTau'$ adnumeraremus. Sed quamquam hoc modo solutionis elegantia sensu mathematico augeretur, tamen ad propositum nostrum praestat, distinctionem illam conservare.

\Article{18}
In applicatione disquisitionis praecedentis ad evolutionem termini secundi expressionis $\Omega$ art.\ 3, spatium secundum inde ab art.\ 6 per $S$ denotatum cum primo identicum est; quae itaque in art.\ 16 erant $\sigma$, $\tau$, $\tau'$, hic erunt $s$, $t$, $0$, si $t$ denotat totam superficiem spatii $s$ a fluido impleti. Quapropter quoties hoc spatium neque partes sensibilis extensionis sed insensibilis crassitiei continet, neque eiusmodi interstitia (fissuras), pars secunda expressionis $\Omega$ fit
\[
=\frac12\pi cc(s\psi0-t\theta0)
\]

Exceptiones itaque adsunt duae:

1) Si spatium $s$ continet partem insensibilis crassitiei, huius superficies duas partes sensibiliter aequales offert, quarum alterutra per $t'$ denotata, crassitieque spatii apud quodvis elementum $\mathrm{d}t'$ indefinite per $\rho$, accedet expressioni praecedenti terminus
\[
\pi cc\int\theta'\rho.\mathrm{d}t'
\]

2) Si spatium $s$ continet cavitatem insensibilis crassitiei, accedet similis terminus, puta $\pi cc\int\theta'\rho.\mathrm{d}t''$, denotante $t''$ alterutram partem superficiei $t$ fissurae contiguam, atque $\rho$ indefinite crassitiem fissurae in quovis puncto.

In evolutione termini tertii expressionis $\Omega$ signum $S$ retinendum erit, ut denotet spatium a vase repletum, sed loco characteristicae $f$ characteristicam $F$ \SourceLeaf{58}ad vim attractivam molecularum vasis relatam substituere oportebit, et perinde loco functionum per characteristicas $\varphi$, $\psi$, $\theta$, $\theta'$ denotatum alias per characteristicas $\Phi$, $\Psi$, $\Theta$, $\Theta'$ denotandas adhibere, quas perinde ab $F$ pendere supponimus ut illas ab $f$. Quae in disquisitione generali erant $\sigma$, $\tau'$, hic manifesto erunt $0$: pro $\tau$ vero hic simpliciter literam $T$ adoptabimus, ut indicet non superficiem totam spatii $S$, sed eam partem, quae fluido contigua est. Hoc pacto pars tertia expressionis $\Omega$ fit, generaliter loquendo,
\[
=\pi cCT\Theta0
\]
exceptis etiam hic duobus casibus, puta

3) Si apud partem sensibilem $T'$ superficiei $T$ fluidum crassitiem insensibilem habet, indefinite per $\rho$ exprimendam, accedet terminus
\[
-\pi cC\int\Theta'\rho.\mathrm{d}T'
\]

4) Si superficies vasis praeter partem $T$ fluido contiguam, offert aliam $T''$ in distantia quidem sed insensibili a fluido positam, accedet, denotante $\rho$ indefinite hac distantiam pro quolibet puncto, terminus
\[
+\pi cC\int\Theta'\rho.\mathrm{d}T''
\]
Superfluum foret, exceptioni primae, quatenus sub tertia non continetur, nec non secundae vel quartae immorari: etiamsi enim aequilibrium fluidi in casibus quibusdam huc referendis, attamen maxime specialibus, locum habere queat, tale aequilibrium nec stabile neque experimentis accessibile esse posset. Contra casus exceptus primus, quatenus sub tertio continetur, utique theoriae essentialis est, verumtamen aliquantisper hic seponetur, ut conditiones aequilibrii, quatenus absque cute fluidi insensibili, vasi adhaerente, consistere potest, explorentur.

Dum itaque omnes has exceptiones seponimus, expressio, cuius valor in statu aequilibrii maximum esse debet, haec erit
\[
-gc\int z\mathrm{d}s+\frac12ccs\psi0-\frac12\pi cct\theta0+\pi cCT\Theta0
\]
et quum in omnibus mutationibus, quas figura fluidi subire potest, spatium $s$ invariatum maneat, expressio sequens
\[
\int z\mathrm{d}s+\frac{\pi c\theta0}{2g}.t-\frac{\pi C\Theta0}{g}.T
\]
in statu aequilibrii \textit{minimum} esse debebit.

\SourceLeaf{59}Iam supra monuimus, $\frac{c\theta0}{g}$ exhibere spatium duarum dimensionum, idemque de $\frac{C\Theta0}{g}$ valet. Statuendo itaque
\[
\frac{\pi c\theta0}{2g}=\alpha\alpha,\qquad\frac{\pi C\Theta0}{2g}=\beta\beta
\]
erunt $\alpha$, $\beta$ lineae constantes a relatione gravitatis ad intensitatem virium, quas partes fluidi a se mutuo et a moleculis vasis patiuntur, pendentes; et si porro partem liberam superficiei fluidi, i.\ e.\ eam, quae vasi non est contigua, per $U$ denotamus. ut habeatur $t=T+U$, minimum esse debet in statu aequilibrii expressio sequens. abhinc per $W$ denotanda:
\[
\int z\mathrm{d}s+(\alpha\alpha-2\beta\beta)T+\alpha\alpha U
\]

\Article{19}
Antequam quae ex hoc theoremate gravissimo sequuntur generaliter et complete evolvamus, operae pretium erit ostendere, quanta facilitate phaenomenon principale tuborum capillarium inde demanet.

Consideremus fluidum in aequilibrio in vase bicrurali, ita ut pars superficiei liberae fluidi sit in primo crure, pars alia in secundo: parietes vasis in confiniis harum partium verticales supponimus. Sit $a$ area sectionis horizontalis internae primi cruris (vel exactius proiectionis horizontalis superficiei liberae fluidi in primo crure), $b$ eiusdem peripheria, denique $ah$ volumen fluidi in hoc crure, pariete verticali deorsum usque ad planum, a quo numerantur distantiae $z$, continuato, sive, quod eodem redit, $h$ altitudo media fluidi supra hoc planum: similia denotentur pro secundo crure per literas $a'$, $b'$, $h'$. Si statum fluidi mutationem infinite parvam subire concipimus, et quidem talem, ut utraque superficiei liberae pars figuram suam servet, variatio partis primae expressionis $W$, puta integralis $\int z\mathrm{d}s$, manifesto erit
\[
=ah\mathrm{d}h+a'h'\mathrm{d}h
\]
variatio ipsius $T$ autem
\[
=b\mathrm{d}h+b'\mathrm{d}h'
\]
denique per hyp $\mathrm{d}U=0$. Hinc colligitur
\[
\mathrm{d}W=ah\mathrm{d}h+a'h'\mathrm{d}h'-(2\beta\beta-\alpha\alpha)(b\mathrm{d}h+b'\mathrm{d}h')
\]
\SourceLeaf{60}Porro quum volumen integrum fluidi invariatum maneat, erit
\[
a\mathrm{d}h+a'\mathrm{d}h'=0
\]
et proin
\[
\mathrm{d}W=\mathrm{d}h[a(h-h')-(2\beta\beta-\alpha\alpha)(b-\frac{ab'}{a'})]
\]
Conditio itaque, ut $W$ in statu aequilibrii sit minimum, perducit ad aequationem, phaenomenon principale tuborum capillarium implicantem
\[
h-h'=(2\beta\beta-\alpha\alpha)(\frac{b}{a}-\frac{b'}{a'})
\]
sponteque patet, huic aequationi revera respondere valorem \textit{minimum} ipsius $W$, quum valor ipsius $\frac{\mathrm{d}\mathrm{d}W}{\mathrm{d}h^2}$ fiat $=a+\frac{aa}{a'}$, i.\ e.\ natura sua positivus

Crus secundum priori largius pronunciatur, si quotiens $\frac{a'}{b'}$ est maior quam $\frac{a}{b}$; fluidum itaque in crure arctiori magis depressum vel magis elevatum erit quam in largiori, prout quadratum $\beta\beta$ minus vel maius est quam $\frac12\alpha\alpha$; et si forte haberetur $\beta\beta=\frac12\alpha\alpha$, altitudo in utroque crure eadem foret. Si crus secundum tam largum est, ut $\frac{b'}{a'}$ negligi possit prae $\frac{b}{a}$, erit proxime
\[
h-h'=(2\beta\beta-\alpha\alpha)\frac{b}{a}
\]
In tubis itaque capillaribus cylindricis fluidi depressio vel elevatio diametro tubo reciproce proportionalis est. Haec omnia tum cum experientia tum cum iis, quae ill.\ \textsc{Laplace} per theoriam stabilire conatus est, conveniunt.

Si vas pluribus cruribus verticalibus inter se communicantibus instructum est, designent $a''$, $b''$, $h''$ pro tertio, $a'''$, $b'''$, $h'''$ pro quarto etc.\ eadem, quae $a$, $b$, $h$ pro primo, eritque etiam
\[
\begin{gathered}
h-h''=(2\beta\beta-\alpha\alpha)(\frac{b}{a}-\frac{b''}{a'})\\
h-h'''=(2\beta\beta-\alpha\alpha)(\frac{b}{a}-\frac{b'''}{a'''})
\end{gathered}
\]
Concinnius hae aequationes ita exhibentur:
\[
\begin{aligned}
h-(2\beta\beta-\alpha\alpha)\frac{b}{a}&=h'-(2\beta\beta-\alpha\alpha)\frac{b'}{a'}\\
&=h''-(2\beta\beta-\alpha\alpha)\frac{b''}{a''}=h'''-(2\beta\beta-\alpha\alpha)\frac{b'''}{a'''}\text{etc.}
\end{aligned}
\]
Quum planum horizontale, a quo altitudines numerantur, arbitrarium sit, patet, si illud ita assumatur, ut sit
\[
h=(2\beta\beta-\alpha\alpha)\frac{b}{a}
\]
\SourceLeaf{61}etiam in reliquis cruribus fore
\[
\begin{gathered}
h'=(2\beta\beta-\alpha\alpha)\frac{b'}{a'},\qquad h''=(2\beta\beta-\alpha\alpha)\frac{b''}{a''},\\
h'''=(2\beta\beta-\alpha\alpha)\frac{b'''}{a'''}\quad\text{etc.}
\end{gathered}
\]
Hocce planum, cuius conceptum infra generalius stabiliemus, vocari potest planum horizontale normale (plan de niveau). Supponendo (si opus est) parietes verticales singulorum crurium usque ad hoc planum productos, $ah$, $a'h'$, $a''h''$ etc.\ expriment, pro $2\beta\beta>\alpha\alpha$, quantitates fluidi in singulis cruribus supra hoc planum elevati, vel pro $2\beta\beta<\alpha\alpha$, quantitates fluidi infra hoc planum in singulis cruribus deficientis: hae itaque quantitates aequales sunt productis ex area constante $2\beta\beta-\alpha\alpha$ in circumferentias $b$, $b'$, $b''$, $b'''$ etc.

\Article{20}
Superest iam, ut e theoremate art.\ 18 indolem figurae aequilibrii determinemus, cuius negotii cardo vertitur in evolutione generali variationis, quam expressio $W$ patitur, dum figura spatii a fluido impleti mutationem quamcunque infinite parvam subit. Sed quum calculus variationum integralium duplicium pro casu, ubi etiam limites tamquam variabiles spectari debent, hactenus parum excultus sit, hanc disquisitionem subtilem paullo profundius petere oportet.

Considerabimus superficiei, quae spatium $s$ a reliquo spatio separat, partem $U$, atque quodvis illius punctum per tres coordinatas $x$, $y$, $z$ determinari supponemus, quarum tertia sit distantia a plano horizontali arbitrario. Spectari itaque poterit $z$ tamquam functio indeterminatarum $x$, $y$, cuius differentialia partialia secundum morem suetum, sed omissis vinculis, per
\[
\frac{\mathrm{d}z}{\mathrm{d}x}.\mathrm{d}x,\qquad\frac{\mathrm{d}z}{\mathrm{d}y}.\mathrm{d}y
\]
denotabimus. In quovis superficiei puncto rectam superficiei normalem et respectu spatii $s$ extrorsum directam concipimus, cosinusque angulorum inter hanc normalem atque rectas axibus coordinatarum $x$, $y$. $z$ parallelas per $\xi$, $\eta$. $\zeta$ denotamus. Hoc pacto erit
\[
\begin{gathered}
\xi\xi+\eta\eta+\zeta\zeta=1\\
\frac{\mathrm{d}z}{\mathrm{d}x}=-\frac{\xi}{\zeta},\qquad\frac{\mathrm{d}z}{\mathrm{d}y}=-\frac{\eta}{\zeta}
\end{gathered}
\]
Limes superficiei $U$ erit linea in se rediens, quam per $P$ denotamus, et dum
